% Lösungen Einheit 2 von 5 – Parallelogramm (zusammenbau v0.1)
\einheitenkopf{Einheit 2 von 5 · Parallelogramm}

\erg{15}{a) Strecke von der oberen Seite senkrecht hinunter auf $g$, rechter Winkel am Fußpunkt \quad b) $A = 9 \cdot 4 = 36$ cm² \quad c) $A = 9 \cdot 5 = 45$ cm² (die schräge Seite zählt nicht) \quad d) $A = 6{,}5 \cdot 4 = 26$ cm² \quad e) $6 \cdot 2 = 12$ cm² und $4 \cdot 3 = 12$ cm² – gleich \quad f) $h = 63 : 9 = 7$ cm \quad g) Höhe senkrecht auf $g$; $A = 6{,}4 \cdot 3{,}5 = 22{,}4$ cm²; ein Dreieck an einer Seite abschneiden und an der anderen anlegen – es entsteht ein Rechteck mit den Seiten $g$ und $h$.}
\erg{16}{a) $u = 2 \cdot (9 + 6) = 30$ cm}
\erg{17}{a) Nils nimmt die schräge Seite statt der Höhe. Richtig: $A = 7 \cdot 4 = 28$ cm² \quad b) Man schneidet an einer Seite ein rechtwinkliges Dreieck ab und legt es an der anderen Seite an. Es entsteht ein Rechteck mit der Länge $g$ und der Breite $h$, also $A = g \cdot h$. \quad c) $A = 45 \cdot 5{,}5 = 247{,}5$ m²; $247{,}5 \cdot 38 = 9\,405$ €}
